% ================================================================
% PROCESSI DI SUPPORTO
% ================================================================

\section{Processi di Supporto}
\label{sec:processi-supporto}

In conformità allo standard ISO/IEC~12207:1995, i processi di
supporto comprendono le attività trasversali necessarie a garantire qualità,
tracciabilità, coerenza e controllo degli artefatti prodotti.

Nel contesto del progetto \textit{Second Brain}, e in accordo con il tailoring
descritto nella Sezione~\ref{sec:tailoring}, il gruppo \textit{Synergo Unit} adotta
i seguenti processi di supporto:

\begin{itemize}
    \item documentazione;
    \item gestione della configurazione;
    \item verifica;
    \item validazione;
    \item accertamento della qualità;
    \item gestione dei cambiamenti;
    \item gestione delle non conformità;
\end{itemize}

\textit{Nota di tailoring sui processi non adottati come autonomi}: lo
standard ISO/IEC~12207:1995 prevede inoltre, tra i processi di supporto, il
Processo di Riesame Congiunto (\textit{Joint Review}, punto 6.6) e il
Processo di Audit (punto 6.7). Il gruppo ha scelto di non normarli come
processi autonomi, ritenendo che il loro intento sia già coperto da pratiche
già normate altrove: la revisione periodica dello stato e dei prodotti di
un'attività (propria del Riesame Congiunto) è svolta tramite la Sprint
Review e la revisione delle Pull Request (Processo di Gestione della
Configurazione, Sezione~\ref{sec:configurazione}); la verifica di conformità
a requisiti, piani e contratto (propria dell'Audit) è svolta dal
Verificatore in sede di Pull Request e nelle revisioni con la proponente.

% ================================================================
% DOCUMENTAZIONE
% ================================================================

\subsection{Processo di Documentazione}
\label{sec:documentazione}

\subsubsection{Descrizione}

Il processo di documentazione definisce regole, strumenti e ciclo di vita per la
produzione, la revisione, l'approvazione e la pubblicazione della documentazione
del gruppo.

Esso si applica a ogni documento gestionale o tecnico, interno o esterno, prodotto
dal gruppo \textit{Synergo Unit}. Il Diario di Bordo\textsubscript{G} segue un template\textsubscript{G} separato
realizzato tramite Canva e non prevede registro delle modifiche.

La documentazione deve essere:
\begin{itemize}
    \item uniforme nella struttura;
    \item coerente con le presenti Norme di Progetto;
    \item verificabile tramite Pull Request\textsubscript{G};
    \item tracciabile rispetto alle modifiche logiche effettuate;
    \item consultabile tramite repository e sito documentale.
\end{itemize}

% ------------------------------------------------

\subsubsection{Tipologia dei Documenti}

\paragraph{Documentazione Gestionale Interna}

\begin{itemize}
    \item \textbf{Norme di Progetto}: definiscono processi, convenzioni, strumenti
    e procedure operative adottate dal gruppo. Sono redatte e mantenute
    dall'Amministratore.

    \item \textbf{Glossario}: raccoglie termini tecnici, acronimi ed espressioni
    ambigue utilizzate nella documentazione. Ogni autore può proporre nuovi
    termini; l'Amministratore ne cura normalizzazione e coerenza.

    \item \textbf{Analisi dello Stato dell'Arte}: documento che raccoglie valutazione comparativa 
    delle tecnologie e motivazione dello stack candidato per il Proof of Concept.

    \item \textbf{Verbali Interni}: documentano riunioni interne, decisioni operative,
    retrospettive, pianificazione e assegnazione dei task\textsubscript{G}.
\end{itemize}

\paragraph{Documentazione Gestionale Esterna}

\begin{itemize}
    \item \textbf{Verbali Esterni}: documentano gli incontri con la proponente e con
    gli stakeholder\textsubscript{G} esterni. Sono redatti dall'Amministratore e
    condivisi con la proponente per conferma e firma.

    \item \textbf{Piano di Progetto}: descrive pianificazione, consuntivazione,
    gestione dei rischi, controllo dei costi e rotazione dei ruoli.

    \item \textbf{Piano di Qualifica}: definisce strategie, metriche e attività di
    verifica e validazione adottate dal gruppo.

    \item \textbf{Diario di Bordo}: materiale gestionale usato per il confronto
        periodico con il committente ed è redatto tramite template Canva.

        Il template del Diario di Bordo deve garantire l'identificazione del materiale
        presentato, riportando in modo chiaro almeno il gruppo, il riferimento al corso
        e la data del Diario di Bordo.

        Ogni pagina, ad eccezione del frontespizio, deve riportare un footer
        contenente almeno:

        \begin{itemize}
            \item numero progressivo del Diario di Bordo;
            \item data nel formato \texttt{AAAA/MM/GG};
            \item numero di pagina e numero totale di pagine;
            \item nome del gruppo;
            \item numero del gruppo.
        \end{itemize}

        La struttura prevista per il footer è:

        \begin{center}
        \texttt{Diario di Bordo N\textdegree{} - AAAA/MM/GG}\\
        \texttt{Synergo Unit - Gruppo 20} \\
        \texttt{numero pagina / numero pagine totali}\\
        \end{center}

    \item \textbf{Sito Web}: punto di accesso alla documentazione pubblicata tramite
    GitHub Pages.
\end{itemize}

\paragraph{Documentazione Tecnica}

\begin{itemize}
    \item \textbf{Analisi dei Requisiti}: documento tecnico esterno che formalizza
    attori, casi d'uso, requisiti, fonti e matrici di tracciabilità.

    \item \textbf{Specifica Tecnica}: documento tecnico esterno, introdotto in fase
    PB, che descrive architettura, decomposizione in componenti, design pattern
    adottati, flusso dei dati e tracciamento componenti-requisiti del prodotto
    finale.

    \item \textbf{Manuale Utente}: documento tecnico esterno, introdotto in fase
    PB, che descrive in linguaggio non tecnico installazione, configurazione e
    utilizzo di ciascuna funzionalità del prodotto, inclusa la risoluzione dei
    problemi più comuni.
\end{itemize}

% ------------------------------------------------

\subsubsection{Strumenti di Documentazione}
\label{sec:strumenti-doc}

La documentazione viene prodotta tramite:

\begin{itemize}
    \item \textbf{\LaTeX{} con Visual Studio Code}: ambiente principale di redazione;

    \item \textbf{LaTeX Workshop}: estensione utilizzata per compilare localmente i
    documenti PDF;

    \item \textbf{Repository GitHub}: archivio ufficiale della documentazione e della
    relativa storia;

    \item \textbf{GitHub Pages}: sito documentale aggiornato periodicamente tramite
    task dedicato;

    \item \textbf{Canva}: strumento utilizzato per il template dei Diari di Bordo;

    \item \textbf{Script Python per Glossario}: strumento utilizzato per applicare
    automaticamente il pedice\textsubscript{G} alla prima occorrenza dei termini
    presenti nel Glossario.
\end{itemize}

Ogni membro lavora esclusivamente su copia locale del repository.

% ------------------------------------------------

\subsubsection{Convenzioni Redazionali}
\label{sec:convenzioni}

Al fine di garantire uniformità e coerenza tra i documenti, il gruppo adotta le
seguenti convenzioni:

\begin{itemize}
    \item \textbf{Nomi dei file}: tutti i nomi devono essere scritti in minuscolo,
    utilizzando il carattere \texttt{\_} come separatore di parola.
    \\
    \textit{Esempio}: \texttt{piano\_di\_progetto.tex}.

    \item \textbf{Nome del file principale}: il file principale \LaTeX{} deve essere
    nominato secondo il nome previsto per il PDF finale.
    \\
    \textit{Esempio}: \texttt{norme\_di\_progetto.tex} produce
    \texttt{norme\_di\_progetto.pdf}.

    \item \textbf{Titoli}: sezioni, sottosezioni e paragrafi devono utilizzare
    Title Case\textsubscript{G} italianizzato, con maiuscola sulle parole semanticamente
    rilevanti e minuscola su articoli, preposizioni e congiunzioni interne.

    \item \textbf{Tabelle lunghe}: le tabelle che possono superare una pagina devono
    essere realizzate tramite \texttt{xltabular} o \texttt{longtable}, in modo da
    consentirne la spezzatura automatica.

    \item \textbf{Pedice\textsubscript{G}}: la prima occorrenza, nell'intero documento, di ogni
    termine presente nel Glossario deve essere marcata con il pedice\textsubscript{G}. Lo script riceve in ingresso la concatenazione dei file
    \LaTeX{} delle sezioni e applica il pedice alla prima occorrenza globale.

    \item \textbf{Link generici}: i collegamenti ipertestuali nel testo devono essere
    prodotti tramite il comando \verb|\link{url}{testo}|.

    \item \textbf{Link documentali in tabelle}: i collegamenti a documenti all'interno
    di tabelle devono essere prodotti tramite il comando
    \verb|\doclink{url}{testo}|, così da preservare leggibilità e andata a capo.

    \item \textbf{Issue GitHub}: i riferimenti a issue devono essere prodotti tramite
    il comando \verb|\ghissue{id}|.

    \item \textbf{Decisioni da verbale interno}: le decisioni interne devono essere
    identificate con codice \texttt{VI-y.x\textsubscript{G}}, dove \texttt{y} è il numero progressivo
    del verbale e \texttt{x} è il numero progressivo della decisione.

    \item \textbf{Decisioni da verbale esterno}: le decisioni esterne devono essere
    identificate con codice \texttt{VE-y.x\textsubscript{G}}, con la stessa logica adottata per i
    verbali interni.
\end{itemize}

% ------------------------------------------------

\subsubsection{Template Documentale}

Tutti i documenti, ad eccezione dei Diari di Bordo, condividono un template
\LaTeX{} comune comprensivo di:

\begin{itemize}
    \item frontespizio;
    \item registro delle modifiche;
    \item indice;
    \item sezioni specifiche del documento;
    \item configurazioni condivise;
    \item loghi e asset comuni.
\end{itemize}

Il template definisce inoltre:
\begin{itemize}
    \item margini;
    \item colori;
    \item profondità dell'indice;
    \item comandi personalizzati;
    \item configurazione dei link;
    \item formattazione delle tabelle;
    \item stile dei paragrafi.
\end{itemize}

% ------------------------------------------------

\subsubsection{Registro delle Modifiche}

Ogni documento \LaTeX{}, ad eccezione dei Diari di Bordo, deve contenere un registro
delle modifiche.

Il registro deve riportare:

\begin{itemize}
    \item \textbf{Versione};
    \item \textbf{Data};
    \item \textbf{Autore};
    \item \textbf{Verificatore};
    \item \textbf{Descrizione}.
\end{itemize}

Il registro delle modifiche è aggiornato dal Redattore.

Le voci del registro devono descrivere modifiche logiche e documentali significative.

% ------------------------------------------------

\subsubsection{Gestione del Glossario}

Il Glossario viene aggiornato in modo incrementale durante la produzione della
documentazione.

Il processo prevede:

\begin{enumerate}
    \item durante la stesura, ogni autore individua eventuali nuovi termini tecnici,
    acronimi o espressioni ambigue;

    \item al termine della stesura, l'autore inserisce i nuovi termini nel file
    \texttt{termini.tex} del Glossario;

    \item l'autore aggiorna il file \texttt{glossary.txt}, utilizzato dallo script
    Python;

    \item l'autore esegue manualmente lo script per applicare il pedice\textsubscript{G}
     alla prima occorrenza dei termini nel documento;

    \item il Verificatore controlla che i termini inseriti siano coerenti con il
    documento e che il pedice\textsubscript{G} sia stato applicato correttamente.
\end{enumerate}

L'Amministratore può intervenire per normalizzare definizioni, duplicati,
maiuscole/minuscole e sinonimi.

% ------------------------------------------------

\subsubsection{Checklist Prima della Pull Request}

Prima di aprire una Pull Request, il Redattore deve verificare che:

\begin{itemize}
    \item il documento compili senza errori bloccanti;
    \item il PDF sia generato correttamente;
    \item il registro delle modifiche sia aggiornato, così come numero di versione e data di ultima modifica;
    \item non siano presenti segnaposto non motivati;
    \item i link siano funzionanti;
    \item i riferimenti interni siano corretti;
    \item le tabelle lunghe siano spezzabili;
    \item il Glossario sia aggiornato, se necessario;
    \item il pedice\textsubscript{G} sia stato applicato correttamente;
    \item il nome del file sia conforme alle convenzioni.
\end{itemize}

% ------------------------------------------------

\subsubsection{Ciclo di Vita Documentale\textsubscript{G}}
\label{sec:ciclo-vita}

Ogni documento attraversa i seguenti stati:

\begin{enumerate}
    \item \textbf{In Stesura}: il Redattore produce o modifica il contenuto su branch
    dedicato;

    \item \textbf{In Verifica}: il Redattore apre una Pull Request; il Verificatore
    controlla il documento e può approvare o richiedere modifiche;

    \item \textbf{Verificato}: il documento ha superato la verifica tramite Pull Request;

    \item \textbf{Approvato}: il documento verificato viene approvato dal Responsabile,
    eventualmente dopo review di gruppo;

    \item \textbf{Rilasciato}: il documento approvato viene pubblicato nel repository
    e reso disponibile tramite sito documentale.
\end{enumerate}

Un documento entra in baseline solo dopo approvazione del Responsabile e review
di gruppo.

% ------------------------------------------------

\subsubsection{Versionamento Documentale}

Il gruppo adotta il \textit{Semantic Versioning\textsubscript{G}} a tre cifre nel
formato:

\begin{center}
    \texttt{vX.Y.Z}
\end{center}

Le versioni vengono incrementate secondo le seguenti regole:

\begin{xltabular}{\textwidth}{
    >{\raggedright\arraybackslash}p{2.5cm}
    >{\raggedright\arraybackslash}p{3cm}
    >{\raggedright\arraybackslash}X
}
    \toprule
    \textbf{Campo} & \textbf{Incremento} & \textbf{Quando Applicarlo} \\
    \midrule
    \endfirsthead

    \toprule
    \textbf{Campo} & \textbf{Incremento} & \textbf{Quando Applicarlo} \\
    \midrule
    \endhead

    \midrule
    \multicolumn{3}{r}{\textit{Continua nella pagina successiva}} \\
    \endfoot

    \bottomrule
    \endlastfoot

    Major
    & \texttt{+1.0.0}
    & Ingresso in una nuova baseline o cambiamento strutturale profondo del documento. \\

    Minor
    & \texttt{+0.1.0}
    & Aggiunta o modifica significativa di contenuti. \\

    Patch
    & \texttt{+0.0.1}
    & Correzioni ortografiche, tipografiche o modifiche non semantiche. \\

\end{xltabular}

L'incremento di versione è scelto dal Redattore, controllato dal Verificatore e
approvato dal Responsabile in caso di rilascio ufficiale.

% ================================================================
% GESTIONE DELLA CONFIGURAZIONE
% ================================================================

\subsection{Processo di Gestione della Configurazione}
\label{sec:configurazione}

\subsubsection{Descrizione}

Il processo di gestione della configurazione garantisce integrità, consistenza,
versionamento e tracciabilità degli elementi configurativi del progetto.

Sono considerati elementi configurativi:

\begin{itemize}
    \item documenti \LaTeX{};
    \item file PDF generati;
    \item template;
    \item asset grafici;
    \item script di automazione;
    \item codice del Proof of Concept, quando presente;
    \item configurazioni del repository;
    \item materiale pubblicato su GitHub Pages.
\end{itemize}

% ------------------------------------------------

\subsubsection{Sistema di Gestione}

Il gruppo utilizza GitHub come sistema principale di gestione della configurazione.

Ogni modifica deve essere tracciata tramite:

\begin{itemize}
    \item issue;
    \item branch dedicato;
    \item commit\textsubscript{G} riferiti alla issue;
    \item Pull Request;
    \item review;
    \item merge\textsubscript{G} controllato.
\end{itemize}

% ------------------------------------------------

\subsubsection{Struttura del Repository}

A partire dallo Sprint 9, la repository ufficiale del progetto assume il nome
\texttt{Second-Brain} e non è più destinata esclusivamente alla documentazione.

Essa contiene:

\begin{itemize}
    \item instruzioni per la CI
    \item documentazione di progetto;
    \item codice del Proof of Concept;
    \item documentazione del MVP;
    \item script di automazione;
    \item configurazioni di supporto;
    \item asset condivisi;
    \item materiali pubblicati tramite GitHub Pages.
\end{itemize}

La struttura principale della repository è:

\begin{verbatim}
Second-Brain/
+-- .github/
+-- docs/
+-- poc/
|   +-- frontend/
|   +-- backend/
+-- mvp
|   +-- frontend/
|   +-- backend/
+-- scripts/
\end{verbatim}

La cartella \texttt{.github/} contiene la checklist e le instruzioni per la CI

La cartella \texttt{docs/} contiene la documentazione di progetto.

La cartella \texttt{poc/} contiene il codice del Proof of Concept, separando
frontend e backend.

La cartella \texttt{mvp/} contine il codice del MVP suddiviso in 
frontend, backend e test.

La cartella \texttt{scripts/} contiene gli script di automazione già adottati dal
gruppo, inclusi quelli a supporto del Glossario e del Piano di Qualifica.

Il sito GitHub Pages pubblica la documentazione e rende disponibile una versione
compressa del PoC per il download.

La versione ufficiale e stabile di ogni elemento configurativo è quella presente
nel branch \texttt{main}\textsubscript{G}.

% ------------------------------------------------

\subsubsection{Gestione delle Issue}

Le issue devono essere create esclusivamente tramite Google Forms.

A partire dallo Sprint 9, il form di creazione delle issue distingue tra attività
documentali e attività non documentali.

\paragraph{Issue documentali}

Per le issue documentali continua a essere applicata la classificazione già
adottata per la documentazione.

Il nome logico della issue documentale deve seguire la forma:

\begin{center}
    \texttt{doc\_<categoria-uso>\_<nome-documento>}
\end{center}

dove:

\begin{itemize}
    \item \texttt{doc} identifica un'attività documentale;
    \item \texttt{categoria-uso} identifica la natura e la destinazione del documento,
    separando i due valori tramite trattino;
    \item \texttt{categoria} può assumere i valori \texttt{gestionale} o
    \texttt{tecnico};
    \item \texttt{uso} può assumere i valori \texttt{interno} o \texttt{esterno};
    \item \texttt{nome-documento} identifica il documento o l'artefatto documentale
    coinvolto, con parole separate da trattini.
\end{itemize}

\textit{Esempi}:

\begin{itemize}
    \item \texttt{doc\_gestionale-interno\_norme-di-progetto};
    \item \texttt{doc\_gestionale-esterno\_piano-di-qualifica};
    \item \texttt{doc\_tecnico-esterno\_analisi-dei-requisiti}.
\end{itemize}

Il sito documentale è considerato un artefatto documentale quando la modifica
riguarda contenuti, collegamenti, indici o documenti pubblicati.

\paragraph{Issue non documentali}

Per le issue non documentali, il form distingue il tipo di attività secondo le
seguenti categorie:

\begin{itemize}
    \item \textbf{feature\textsubscript{G}}: nuova funzionalità del prodotto o del PoC;
    \item \textbf{bug\textsubscript{G}}: comportamento non corretto rispetto all'atteso;
    \item \textbf{refactor\textsubscript{G}}: miglioramento interno senza modifica del comportamento
    osservabile;
    \item \textbf{test}: attività relativa ai test, prevista per le attività successive
    alla RTB.
\end{itemize}

La posizione deve essere scelta tra:

\begin{itemize}
    \item \textbf{poc};
    \item \textbf{frontend};
    \item \textbf{backend};
    \item \textbf{infra};
    \item \textbf{test}.
\end{itemize}

Il nome logico della issue non documentale deve seguire la forma:

\begin{center}
    \texttt{<tipo>\_<posizione>\_<descrizione-breve>}
\end{center}

La separazione tra i campi principali avviene tramite il carattere
\texttt{\_}. La descrizione breve deve essere scritta in minuscolo, senza accenti
e con parole separate da trattini.

\textit{Esempi}:

\begin{itemize}
    \item \texttt{feature\_frontend\_css};
    \item \texttt{bug\_frontend\_css};
    \item \texttt{refactor\_infra\_aggiornamento-token};
    \item \texttt{test\_backend\_test-chiamata-ai}.
\end{itemize}

\paragraph{Campi obbligatori}

Ogni issue deve contenere:

\begin{itemize}
    \item tipo o categoria;
    \item posizione, uso o ambito;
    \item assegnatario;
    \item ruolo;
    \item priorità;
    \item scadenza;
    \item tempo previsto\textsubscript{G};
    \item milestone RTB o PB;
    \item descrizione dell'attività.
\end{itemize}

Le modifiche strutturali o configurative del sito documentale sono invece
classificate come attività infrastrutturali con posizione \texttt{infra}.

% ------------------------------------------------

\subsubsection{Strategia di Branching}

Il gruppo adotta un modello \textit{trunk-based development\textsubscript{G}}.

Il branch \texttt{main} è l'unico branch stabile. Ogni modifica deve essere sviluppata su branch dedicato, creato a partire dalla
sezione \textit{Development} della issue GitHub associata.

A partire dallo Sprint 9, con decorrenza dal 2026-06-02, il gruppo distingue tra
attività documentali e attività non documentali.

Per tutte le nuove convenzioni di naming, il carattere \texttt{\_} separa i campi
principali, mentre il carattere \texttt{-} separa le parole interne ai singoli
campi composti.

\paragraph{Branch documentali}

Per le attività documentali, la convenzione di naming dei branch è:

\begin{center}
    \texttt{<id>\_doc\_<categoria-uso>\_<nome-documento>}
\end{center}

dove:

\begin{itemize}
    \item \texttt{id} è l'identificativo della issue;
    \item \texttt{doc} identifica un'attività documentale;
    \item \texttt{categoria-uso} indica natura e destinazione del documento,
    separando i due valori tramite trattino;
    \item \texttt{categoria} può assumere i valori \texttt{gestionale} o
    \texttt{tecnico};
    \item \texttt{uso} può assumere i valori \texttt{interno} o \texttt{esterno};
    \item \texttt{nome-documento} identifica il documento o l'artefatto documentale
    modificato, con parole separate da trattini.
\end{itemize}

\textit{Esempi}:

\begin{itemize}
    \item \texttt{146\_doc\_gestionale-interno\_norme-di-progetto};
    \item \texttt{147\_doc\_gestionale-esterno\_piano-di-qualifica};
    \item \texttt{148\_doc\_tecnico-esterno\_analisi-dei-requisiti}.
\end{itemize}

\paragraph{Branch non documentali}

Per le attività non documentali, la convenzione di naming dei branch è:

\begin{center}
    \texttt{<id>\_<tipo>\_<posizione>\_<descrizione-breve>}
\end{center}

dove:

\begin{itemize}
    \item \texttt{id} è l'identificativo della issue;
    \item \texttt{tipo} indica la natura dell'attività e può assumere i valori
    \texttt{feature}, \texttt{bug}, \texttt{refactor} o \texttt{test};
    \item \texttt{posizione} indica l'area coinvolta e può assumere i valori
    \texttt{poc}, \texttt{frontend}, \texttt{backend}, \texttt{infra} o \texttt{test};
    \item \texttt{descrizione-breve} sintetizza l'attività con parole minuscole,
    senza accenti e separate da trattini.
\end{itemize}

La separazione tra i campi principali avviene tramite il carattere
\texttt{\_}; la separazione tra le parole interne alla descrizione breve avviene
tramite trattino.

\textit{Esempi}:

\begin{itemize}
    \item \texttt{142\_feature\_frontend\_css};
    \item \texttt{143\_bug\_frontend\_css};
    \item \texttt{144\_refactor\_infra\_aggiornamento-token};
    \item \texttt{145\_test\_backend\_test-chiamata-ai}.
\end{itemize}

Le issue e i branch già aperti prima dell'adozione della presente convenzione
mantengono la nomenclatura originaria fino alla loro chiusura. Le nuove attività
aperte a partire dallo Sprint 9 devono seguire le convenzioni aggiornate.

% ------------------------------------------------

\subsubsection{Commit}

Ogni commit deve seguire una convenzione ispirata ai Conventional Commits\textsubscript{G},
adattata al workflow del gruppo.

La forma prescritta è:

\begin{center}
    \texttt{<tipo>(<posizione>): <messaggio> (\#id)}
\end{center}

dove:

\begin{itemize}
    \item \texttt{tipo} identifica la natura della modifica;
    \item \texttt{posizione} indica l'area o il documento coinvolto;
    \item \texttt{messaggio} descrive sinteticamente la modifica;
    \item \texttt{id} è l'identificativo della issue collegata.
\end{itemize}

I tipi ammessi sono:

\begin{xltabular}{\textwidth}{
    >{\raggedright\arraybackslash}p{3cm}
    >{\raggedright\arraybackslash}p{3cm}
    >{\raggedright\arraybackslash}X
}
    \toprule
    \textbf{Tipo Task} & \textbf{Tipo Commit} & \textbf{Significato} \\
    \midrule
    \endfirsthead

    \toprule
    \textbf{Tipo Task} & \textbf{Tipo Commit} & \textbf{Significato} \\
    \midrule
    \endhead

    \midrule
    \multicolumn{3}{r}{\textit{Continua nella pagina successiva}} \\
    \endfoot

    \bottomrule
    \endlastfoot

    feature
    & \texttt{feat}
    & Nuova funzionalità del prodotto o del PoC. \\

    bug
    & \texttt{fix}
    & Correzione di un comportamento non corretto. \\

    doc
    & \texttt{docs}
    & Modifica alla documentazione. \\

    refactor
    & \texttt{refactor}
    & Miglioramento interno senza modifica del comportamento osservabile. \\

    test
    & \texttt{test}
    & Aggiunta o modifica di test. \\

\end{xltabular}

Il tipo del commit deve essere coerente con il tipo della issue associata (per le attività documentali, il tipo di commit è sempre \texttt{docs}.).

Il messaggio deve essere scritto in italiano e in forma descrittiva.

Ogni commit deve riferirsi a una sola issue. Le Pull Request possono includere più
issue solo quando le attività sono strettamente collaborative o tecnicamente
indivisibili.

\paragraph{Scope Documentali}

Per i commit documentali devono essere utilizzati i seguenti scope\textsubscript{G}:

\begin{itemize}
    \item \texttt{ndp}: Norme di Progetto;
    \item \texttt{pdq}: Piano di Qualifica;
    \item \texttt{pdp}: Piano di Progetto;
    \item \texttt{adr}: Analisi dei Requisiti;
    \item \texttt{sa}: Analisi dello Stato dell'Arte;
    \item \texttt{glo}: Glossario;
    \item \texttt{viN}: Verbale interno numero \texttt{N};
    \item \texttt{veN}: Verbale esterno numero \texttt{N};
    \item \texttt{ddbN}: Diario di Bordo numero \texttt{N};
    \item \texttt{sito}: sito documentale.
    \item \texttt{st}: Specifica tecnica.
    \item \texttt{mu}: Manuale utente.
\end{itemize}

Lo scope \texttt{adr} indica l'Analisi dei Requisiti nel contesto dei commit
documentali e non deve essere confuso con gli Architecture Decision Records.

\paragraph{Esempi}

\begin{itemize}
    \item \texttt{feat(frontend): aggiunto css (\#142)};
    \item \texttt{fix(frontend): sistemato problema padding css (\#143)};
    \item \texttt{docs(ndp): aggiunta sezione 2 (\#144)};
    \item \texttt{refactor(infra): aggiornato token (\#145)};
    \item \texttt{test(backend): aggiunto test chiamata ai (\#146)}.
\end{itemize}

% ------------------------------------------------

\subsubsection{Pull Request}

Ogni modifica deve essere integrata tramite Pull Request.

La Pull Request deve:

\begin{itemize}
    \item essere collegata alla issue corrispondente;
    \item descrivere sinteticamente le modifiche effettuate;
    \item essere aperta preferibilmente entro la domenica precedente la riunione del
    martedì;
    \item ricevere almeno una review positiva;
    \item non essere approvata dall'autore della modifica.
\end{itemize}

In caso di modifiche sostanziali richieste dal Verificatore, il Redattore può
segnare la Pull Request come \textit{draft} fino al completamento delle correzioni.

In fase PB, la \textit{Pull Request Policy} viene rafforzata:
\begin{itemize}
    \item È obbligatoria l'approvazione di almeno \textbf{1 reviewer} diverso dall'autore del codice. 
    \item Il reviewer ha l'esplicito compito di verificare non solo la logica funzionale, ma anche il rigoroso rispetto delle regole di dipendenza tra gli strati architetturali.
\end{itemize}

% ------------------------------------------------

\paragraph{Pull Request per branch di codice}

Per i branch di codice la prassi in ordine è:

\begin{itemize}
    \item creazione del branch per lo sviluppo,
    \item creare la PR in draft verso il main,
    \item prima di ogni push bisogna fare i test in locale,
    \item se si passano i test in locale, si fa il push e la PR in draft fa partire la CI,
    \item terminato lo sviluppo e compilata la checklist, si mette in \textit{ready for review},
    \item il verificatore fa la review e approva solo se la pipeline risulta interamente verde.
\end{itemize}


% ------------------------------------------------

\subsubsection{Branch Protection Rules}

Il branch \texttt{main} è protetto tramite branch protection rules\textsubscript{G}.

Le regole attive prevedono:

\begin{itemize}
    \item obbligo di almeno una review positiva;
    \item blocco dei force push\textsubscript{G};
    \item restrizione dell'eliminazione del branch protetto.
\end{itemize}

% ------------------------------------------------

\subsubsection{Merge}

Il merge su \texttt{main} può essere eseguito solo dopo approvazione della Pull
Request da parte del Verificatore.

Il merge viene eseguito dal Redattore, non dal Verificatore.

Il gruppo utilizza la modalità standard di merge proposta da GitHub tramite Pull
Request.

Dopo il merge:

\begin{itemize}
    \item la issue associata viene chiusa automaticamente;
    \item il branch di lavoro viene eliminato;
    \item il Redattore aggiorna il tempo effettivo\textsubscript{G} nel foglio di consuntivazione.
\end{itemize}

% ------------------------------------------------

\subsubsection{Generazione e Pubblicazione dei PDF}

I PDF vengono generati localmente tramite LaTeX Workshop.

Il sito GitHub Pages non viene aggiornato automaticamente a ogni merge su
\texttt{main}; l'aggiornamento del sito documentale viene gestito tramite task
periodica dedicata.

% ------------------------------------------------

\subsubsection{Gestione del Sito Documentale}

Il sito documentale pubblicato tramite GitHub Pages costituisce un elemento
configurativo del progetto e viene aggiornato mediante task dedicate.

Le attività relative al sito documentale sono distinte in due categorie:

\begin{itemize}
\item \textbf{Aggiornamenti documentali}: comprendono la pubblicazione o
sostituzione di documenti PDF, l'aggiornamento degli indici, dei collegamenti,
delle pagine informative e di ogni altro contenuto documentale reso disponibile
sul sito.

Tali attività sono considerate documentali e devono seguire le convenzioni
previste per documenti, issue e branch documentali.

Per queste attività deve essere utilizzato lo scope
\texttt{sito} nei commit.

\textit{Esempi}:

\begin{itemize}
    \item branch:
    \texttt{215\_doc\_gestionale-esterno\_sito-documentale};

    \item commit:
    \texttt{docs(sito): aggiornato indice documenti (\#215)}.
\end{itemize}

\item \textbf{Aggiornamenti infrastrutturali}: comprendono modifiche alla
struttura del sito, alla navigazione, al tema grafico, agli script di build,
ai workflow GitHub Pages e alle configurazioni tecniche necessarie alla
pubblicazione.

Tali attività sono considerate non documentali e devono essere classificate
come attività di tipo \texttt{feature}, \texttt{bug} o \texttt{refactor}
con posizione \texttt{infra}.

\textit{Esempi}:

\begin{itemize}
    \item branch:
    \texttt{216\_feature\_infra\_navbar-sito};

    \item commit:
    \texttt{feat(infra): aggiunta navbar sito (\#216)}.
\end{itemize}
\end{itemize}

L'aggiornamento periodico del sito documentale finalizzato alla pubblicazione
dei documenti approvati è considerato attività documentale.

% ================================================================
% VERIFICA
% ================================================================

\subsection{Processo di Verifica}
\label{sec:verifica}

\subsubsection{Descrizione}

Il processo di verifica assicura che ogni artefatto prodotto sia conforme alle
presenti norme e agli obiettivi di qualità definiti nel Piano di Qualifica.

La verifica risponde alla domanda:

\begin{center}
    \textit{Il prodotto o artefatto è stato realizzato correttamente rispetto alle norme stabilite?}
\end{center}

In fase RTB e durante l'implementazione del Proof of Concept, la verifica riguarda:

\begin{itemize}
    \item documentazione;
    \item coerenza dei riferimenti;
    \item correttezza delle matrici di tracciabilità;
    \item rispetto delle convenzioni;
    \item adeguatezza del materiale relativo al Proof of Concept;
    \item codice prodotto per il Proof of Concept;
    \item configurazioni necessarie all'esecuzione locale del PoC.
\end{itemize}

In fase PB, con l'estensione dello sviluppo al Minimum Viable Product, l'ambito della verifica si allarga a:

\begin{itemize}
    \item codice di produzione del MVP, incluso il rispetto delle convenzioni di nomenclatura (Sezione~\ref{sec:sviluppo});
    \item esito dell'analisi statica (linting e formattazione) e dei test automatici (unitari e di integrazione);
    \item rispetto delle soglie minime di code coverage definite (90\% Backend, 85\% Frontend);
    \item esito della pipeline di Continuous Integration su ogni Pull Request;
    \item Specifica Tecnica e Manuale Utente, introdotti in questa fase;
    \item coerenza dei diagrammi UML (classi, sequenza, package) con il codice effettivamente prodotto.
\end{itemize}

% ------------------------------------------------

\subsubsection{Responsabilità}

\begin{itemize}
    \item Il Verificatore deve essere sempre un membro diverso dal Redattore.
    \item Nessun membro può verificare e approvare il proprio lavoro.
    \item Il Verificatore non modifica direttamente l'artefatto verificato.
    \item Il Verificatore segnala le modifiche richieste tramite commenti o richieste
    di cambiamento nella Pull Request.
    \item Il Responsabile interviene nell'approvazione finale, non nella verifica
    puntuale dell'artefatto.
\end{itemize}

% ------------------------------------------------

\subsubsection{Attività di Verifica}

Per i task documentali, il Verificatore deve controllare:

\begin{itemize}
    \item correttezza linguistica;
    \item coerenza contenutistica;
    \item rispetto del template;
    \item correttezza del registro delle modifiche;
    \item correttezza dei riferimenti interni;
    \item funzionamento dei link;
    \item presenza del pedice\textsubscript{G} dove necessario;
    \item coerenza con il Glossario;
    \item conformità alle convenzioni di naming;
    \item correttezza del versionamento;
    \item compilazione corretta del documento PDF;
    \item assenza di segnaposto non motivati.
\end{itemize}

Se la modifica coinvolge documenti collegati, il Verificatore deve controllare
anche i riferimenti coinvolti.

In particolare, modifiche ai requisiti devono comportare il controllo di:

\begin{itemize}
    \item Analisi dei Requisiti;
    \item matrici di tracciabilità;
    \item Glossario, se necessario;
    \item eventuali riferimenti nel Piano di Qualifica.
\end{itemize}

% ------------------------------------------------

\subsubsection{Verifica del Codice del Proof of Concept}

Durante la fase RTB, la verifica del codice riguarda esclusivamente il codice
prodotto per il Proof of Concept.

Per le Pull Request di implementazione, il Verificatore deve controllare:

\begin{itemize}
    \item coerenza della modifica rispetto alla issue;
    \item correttezza della convenzione di branch e commit;
    \item avvio locale del PoC;
    \item assenza di errori evidenti in console;
    \item assenza di credenziali, token o file \texttt{.env} nel repository;
    \item aggiornamento del README\textsubscript{G} se la modifica cambia la procedura di avvio;
    \item descrizione della prova effettuata per modifiche backend o API\textsubscript{G};
    \item coerenza con requisiti e casi d'uso collegati, se disponibili.
\end{itemize}

La verifica del codice PoC non prevede test automatici obbligatori in fase RTB.
Tali attività saranno normate quando il gruppo introdurrà il processo di testing
nelle fasi successive.

% ------------------------------------------------

\subsubsection{Verifica del Codice (Analisi Statica e Testing)}
In fase PB, il codice prodotto è soggetto ad analisi statica e testing automatico prima di poter essere integrato:
\begin{itemize}
    \item \textbf{Analisi Statica (Linting/Formattazione):} \texttt{ruff} per il backend (Python) e la combinazione \texttt{eslint + prettier} per il frontend (TypeScript/Vue).
    \item \textbf{Testing Automatico:} \texttt{pytest} per il backend e \texttt{vitest} per il frontend.
    \item \textbf{Code Coverage:} sono fissate soglie minime di copertura dei test differenziate per componente: \textbf{90\%} per il Backend (\texttt{pytest-cov}), \textbf{85\%} di linee per il Frontend (\texttt{vitest}); le soglie sono verificate automaticamente in pipeline e fanno fallire la build se non rispettate.
\end{itemize}
Tali controlli sono automatizzati tramite pipeline di Continuous Integration. 
Per ottimizzare i tempi, la pipeline di Continuous Integration è configurata per attivarsi esclusivamente in seguito a modifiche all'interno della directory del codice sorgente (\texttt{mvp/}), ignorando i commit puramente documentali.

% ------------------------------------------------

\subsubsection{Esito della Verifica}

La verifica può avere:

\begin{itemize}
    \item \textbf{Esito positivo}: il Verificatore approva la Pull Request;
    \item \textbf{Esito negativo}: il Verificatore richiede modifiche tramite commenti
    o richiesta di cambiamento.
\end{itemize}

Una richiesta di modifica può bloccare il merge finché il Redattore non ha risolto
le osservazioni.

% ------------------------------------------------

\subsubsection{Definition of Done}
\label{sec:dod}

Un task è considerato \textit{Done} solo se soddisfa i criteri comuni e, in base
alla natura dell'attività, i criteri specifici per documentazione o implementazione.

\paragraph{Criteri Comuni}

\begin{itemize}
    \item il task è stato creato tramite Google Form\textsubscript{G};
    \item il task è associato a una issue GitHub;
    \item la issue contiene tipo, posizione, assegnatario, ruolo, priorità, scadenza
    e tempo previsto;
    \item il task è stato sviluppato su branch dedicato;
    \item i commit rispettano la convenzione definita nelle presenti norme;
    \item la Pull Request è collegata alla issue corrispondente;
    \item la Pull Request ha ricevuto almeno una review positiva;
    \item il Verificatore è diverso dall'autore della modifica;
    \item il Redattore ha eseguito il merge;
    \item la issue è stata chiusa;
    \item il tempo effettivo è stato registrato.
\end{itemize}

\paragraph{Criteri per Task Documentali}

\begin{itemize}
    \item il contenuto è completo rispetto allo scope della issue;
    \item il registro delle modifiche è stato aggiornato, se necessario;
    \item il documento compila correttamente;
    \item il PDF è stato generato senza errori bloccanti;
    \item i link e i riferimenti interni sono corretti;
    \item il Glossario è stato aggiornato, se necessario;
    \item il pedice\textsubscript{G} è stato applicato correttamente.
\end{itemize}

\paragraph{Criteri per Task di Implementazione del PoC}

\begin{itemize}
    \item il codice è coerente con lo scope della issue;
    \item il PoC si avvia localmente;
    \item non sono presenti errori evidenti in console;
    \item non sono presenti credenziali, token, file \texttt{.env} o dati sensibili
    nel repository;
    \item il README è aggiornato se la modifica cambia la procedura di avvio;
    \item l'Analisi dei Requisiti è aggiornata se emerge una modifica funzionale;
    \item l'Analisi dello Stato dell'Arte è aggiornata se cambia una tecnologia;
    \item il Piano di Qualifica è aggiornato se cambiano metriche o modalità di
    misurazione\textsubscript{G};
    \item il Glossario è aggiornato se vengono introdotti nuovi termini tecnici;
    \item per le issue di tipo \texttt{bug}, la Pull Request descrive la causa del
    problema e la correzione effettuata;
    \item per le issue di tipo \texttt{feature}, la Pull Request indica il requisito
    o il caso d'uso collegato, se disponibile.
\end{itemize}

\paragraph{Criteri per Task di Sviluppo (Fase PB)}

Oltre ai criteri comuni, un task di sviluppo del prodotto finale si considera completato solo se:
\begin{itemize}
    \item il codice compila e non sono presenti errori rilevati dai tool di analisi statica (linting);
    \item i test di unità sono stati scritti e risultano eseguiti con successo;
    \item la \textit{code coverage} del modulo modificato o aggiunto rispetta la soglia minima definita per il componente (90\% Backend, 85\% Frontend; Sezione~\ref{sec:verifica});
    \item la Pull Request è approvata da almeno un revisore, che ha controllato il rispetto dei layer architetturali;
    \item in caso di modifiche a interfacce condivise tra componenti, la modifica è stata supervisionata dall'architetto o referente di coerenza;
    \item la pipeline di Continuous Integration è passata con successo.
\end{itemize}

% ================================================================
% VALIDAZIONE
% ================================================================

\subsection{Processo di Validazione}
\label{sec:validazione}

\subsubsection{Descrizione}

Il processo di validazione ha lo scopo di accertare che quanto prodotto risponda
alle esigenze della proponente e agli obiettivi del capitolato.

La validazione risponde alla domanda:

\begin{center}
    \textit{Stiamo costruendo il prodotto giusto?}
\end{center}

In fase RTB, la validazione riguarda principalmente:

\begin{itemize}
    \item requisiti;
    \item casi d'uso;
    \item scelte tecnologiche;
    \item perimetro del Proof of Concept.
\end{itemize}

In fase PB, con il completamento del Minimum Viable Product, la validazione riguarda:

\begin{itemize}
    \item soddisfacimento della totalità dei requisiti obbligatori (Sezione~\ref{sec:sviluppo});
    \item corrispondenza tra funzionalità implementate e casi d'uso dell'Analisi dei Requisiti, tramite la Mappatura Test di Sistema - Casi d'Uso;
    \item architettura e design pattern effettivamente adottati, rispetto a quanto documentato nella Specifica Tecnica;
    \item usabilità e chiarezza del Manuale Utente per un utente privo di competenze tecniche pregresse.
\end{itemize}

% ------------------------------------------------

\subsubsection{Validazione con la Proponente}

Ogni incontro di validazione con la proponente deve produrre un verbale esterno.

La issue relativa alla redazione del verbale esterno rimane aperta fino a quando
il verbale non è stato firmato o confermato dalla proponente via email.

Se la proponente richiede modifiche rilevanti, queste devono essere:

\begin{itemize}
    \item riportate nel verbale esterno;
    \item tradotte in issue dedicate;
    \item implementate tramite Pull Request;
    \item verificate dal Verificatore;
    \item integrate nei documenti coinvolti.
\end{itemize}

Le modifiche ai requisiti derivanti da feedback esterno devono aggiornare:

\begin{itemize}
    \item Analisi dei Requisiti;
    \item matrici di tracciabilità;
    \item Glossario, se necessario;
    \item Piano di Qualifica, se impattano test, metriche o criteri di accettazione.
\end{itemize}

% ================================================================
% ACCERTAMENTO DELLA QUALITÀ
% ================================================================

\subsection{Processo di Accertamento della Qualità}
\label{sec:qualita}

\subsubsection{Descrizione}

Il processo di accertamento della qualità ha lo scopo di monitorare e migliorare
la qualità dei processi e degli artefatti prodotti.

A differenza della verifica, che controlla puntualmente un artefatto, l'accertamento
della qualità monitora l'efficacia e l'efficienza del way of working.

% ------------------------------------------------

\subsubsection{Modello di Miglioramento}

Il gruppo adotta il ciclo PDCA\textsubscript{G} come riferimento per il miglioramento
continuo:

\begin{itemize}
    \item \textbf{Plan}: definizione di obiettivi, metriche e procedure;
    \item \textbf{Do}: applicazione dei processi definiti;
    \item \textbf{Check}: misurazione dei risultati e analisi degli scostamenti;
    \item \textbf{Act}: introduzione di azioni correttive e aggiornamento del way of working.
\end{itemize}

% ------------------------------------------------

\subsubsection{Metriche e Piano di Qualifica}

Il Piano di Qualifica definisce le metriche, le soglie e il cruscotto di qualità\textsubscript{G}
utilizzati dal gruppo per monitorare qualità di processo e qualità di prodotto.

Le Norme di Progetto disciplinano il processo di raccolta, aggiornamento e gestione
delle metriche; formule, soglie, valori rilevati e rappresentazioni grafiche sono
riportati nel \link{https://synergo-unit-unipd.github.io/Second-Brain/docs/PB/doc_gestionali/doc_esterni/piano_di_qualifica/piano_di_qualifica.pdf}{Piano di Qualifica}.

L'Amministratore assegnato alla relativa task raccoglie i dati necessari, aggiorna
il file di tracciamento e mantiene il cruscotto delle metriche.

Le metriche vengono rilevate e aggiornate al termine di ogni Sprint, integrando i dati raccolti nello
sprint concluso e aggiornando le rappresentazioni tabellari e grafiche presenti
nel Piano di Qualifica.

Le metriche relative al codice vengono calcolate solo sul codice mergiato nel
branch \texttt{main}.

Se una metrica risulta fuori soglia, il gruppo ne discute durante la retrospective.
L'esito della discussione viene verbalizzato e, qualora richieda una modifica
operativa, viene creata una issue dedicata.

\paragraph{Metriche di qualità di processo}

Le metriche di qualità di processo misurano l'efficienza, la stabilità e la
controllabilità del way of working adottato dal gruppo.

Il Piano di Qualifica definisce le seguenti metriche di processo:

\begin{xltabular}{\textwidth}{
    >{\raggedright\arraybackslash}p{2.5cm}
    >{\raggedright\arraybackslash}X
    >{\raggedright\arraybackslash}p{6cm}
}
    \toprule
    \textbf{Codice} & \textbf{Metrica} & \textbf{Finalità} \\
    \midrule
    \endfirsthead

    \toprule
    \textbf{Codice} & \textbf{Metrica} & \textbf{Finalità} \\
    \midrule
    \endhead

    \midrule
    \multicolumn{3}{r}{\textit{Continua nella pagina successiva}} \\
    \endfoot

    \bottomrule
    \endlastfoot

    MPC01
    & Schedule Performance Index\textsubscript{G} (SPI\textsubscript{G})
    & Misurare l'efficienza nel rispetto dei tempi pianificati. \\

    MPC02
    & Cost Performance Index\textsubscript{G} (CPI\textsubscript{G})
    & Misurare l'efficienza nell'utilizzo del budget allocato. \\

    MPC03
    & Requirements Stability Index\textsubscript{G} (RSI\textsubscript{G})
    & Monitorare la stabilità dei requisiti e le variazioni introdotte nell'Analisi dei Requisiti. \\

    MPC04
    & Time Efficiency\textsubscript{G}
    & Misurare la percentuale di ore produttive rispetto al totale delle ore spese. \\

    MPC05
    & Task Completion Rate\textsubscript{G}
    & Misurare la percentuale di issue chiuse rispetto a quelle pianificate. \\

    MPC06
    & PR Resolution Time
    & Misurare il tempo medio necessario per la revisione e il merge di una Pull Request. \\

    MPC07
    & Review Coverage\textsubscript{G}
    & Misurare la percentuale di Pull Request approvate da un revisore diverso dall'autore. \\

    MPC08
    & Quality Metrics Compliance Rate\textsubscript{G}
    & Misurare la percentuale di metriche di qualità il cui valore rilevato è almeno accettabile rispetto alle soglie definite. \\

\end{xltabular}

\paragraph{Metriche di qualità di prodotto}

Le metriche di qualità di prodotto misurano la qualità degli artefatti prodotti,
sia documentali sia software.

Il Piano di Qualifica definisce le seguenti metriche di prodotto:

\begin{xltabular}{\textwidth}{
    >{\raggedright\arraybackslash}p{2.5cm}
    >{\raggedright\arraybackslash}X
    >{\raggedright\arraybackslash}p{6cm}
}
    \toprule
    \textbf{Codice} & \textbf{Metrica} & \textbf{Finalità} \\
    \midrule
    \endfirsthead

    \toprule
    \textbf{Codice} & \textbf{Metrica} & \textbf{Finalità} \\
    \midrule
    \endhead

    \midrule
    \multicolumn{3}{r}{\textit{Continua nella pagina successiva}} \\
    \endfoot

    \bottomrule
    \endlastfoot

    MPD01
    & Soddisfacimento Requisiti Obbligatori
    & Misurare la percentuale di requisiti obbligatori implementati. \\

    MPD02
    & Comment Density\textsubscript{G}
    & Misurare il rapporto tra righe di commento e righe di codice totali. \\

    MPD03
    & Browser Compatibility\textsubscript{G}
    & Misurare il numero di browser target su cui il sistema è testato. \\

    MPD04
    & Errori Ortografici
    & Misurare il numero di errori ortografici rilevati per pagina di documento. \\

    MPD05
    & Complessità Ciclomatica\textsubscript{G}
    & Misurare la complessità dei flussi di controllo delle funzioni del codice frontend e backend. \\

    MPD06
    & Tempo Iniziale di Caricamento del Programma\textsubscript{G}
    & Misurare il tempo necessario affinché l'interfaccia principale del PoC risulti disponibile e utilizzabile. \\

    MPD07
    & Indice di Gulpease\textsubscript{G}
    & Misurare la leggibilità dei documenti in lingua italiana. \\

\end{xltabular}

Le formule, gli strumenti di dettaglio, le soglie e le modalità operative di
misurazione sono riportati nel Piano di Qualifica.

% ================================================================
% GESTIONE DEI CAMBIAMENTI
% ================================================================

\subsection{Processo di Gestione dei Cambiamenti}
\label{sec:cambiamenti}

\textit{Nota di tailoring}: lo standard ISO/IEC~12207:1995 non prevede un
processo di supporto autonomo per la gestione dei cambiamenti: il controllo
delle modifiche vi è trattato come parte del Processo di Gestione della
Configurazione (Sezione~\ref{sec:configurazione}). Il gruppo ha scelto di
scorporarlo in un processo a sé, estensione deliberata oltre lo standard, per
dare visibilità e tracciabilità autonome anche ai cambiamenti che non
riguardano direttamente elementi in configurazione (es. cambiamenti di
processo o al way of working) e che la sola Gestione della Configurazione non
coprirebbe.

\subsubsection{Descrizione}

Il processo di gestione dei cambiamenti disciplina la valutazione, l'approvazione,
l'implementazione e la verifica delle modifiche rilevanti a documenti, requisiti,
tecnologie e processi.

Ogni cambiamento rilevante deve essere tracciato.

% ------------------------------------------------

\subsubsection{Tipologie di Cambiamento}

I cambiamenti vengono distinti in:

\begin{itemize}
    \item \textbf{documentali}: modifiche a struttura o contenuto dei documenti;
    \item \textbf{di requisito}: modifiche a requisiti, casi d'uso o matrici di tracciabilità;
    \item \textbf{tecnologici}: modifiche allo stack o alle tecnologie candidate;
    \item \textbf{di processo}: modifiche al way of working o alle presenti norme.
\end{itemize}

% ------------------------------------------------

\subsubsection{Procedura}

Ogni cambiamento segue la seguente procedura:

\begin{enumerate}
    \item \textbf{Richiesta}: il cambiamento emerge da verbale interno, verbale esterno,
      retrospettiva, feedback della proponente o revisione esterna;

    \item \textbf{Valutazione Di Impatto}: il gruppo valuta impatto su documenti,
    requisiti, tempi, costi, tecnologie o processi;

    \item \textbf{Approvazione}: il cambiamento viene approvato dal ruolo competente;

    \item \textbf{Implementazione}: viene creata una issue e il cambiamento viene
    implementato tramite branch e Pull Request;

    \item \textbf{Verifica}: il Verificatore controlla la conformità della modifica;

    \item \textbf{Chiusura}: la issue viene chiusa dopo merge e aggiornamento del
    tempo effettivo.
\end{enumerate}

% ------------------------------------------------

\subsubsection{Responsabilità di Approvazione}

\begin{itemize}
    \item I cambiamenti documentali ordinari sono approvati tramite verifica della
    Pull Request.

    \item I cambiamenti ai requisiti sono approvati dall'Analista e dal Responsabile;
    se incidono sul perimetro funzionale, devono essere validati con la proponente.

    \item I cambiamenti tecnologici sono approvati dal Progettista e dal Responsabile;
    se rilevanti, devono essere discussi con la proponente.

    \item I cambiamenti alle Norme di Progetto sono approvati dall'Amministratore e
    dal Responsabile.
\end{itemize}

% ================================================================
% GESTIONE DELLE NON CONFORMITÀ
% ================================================================

\subsection{Processo di Gestione delle Non Conformità}
\label{sec:non-conformita}

\textit{Nota di tailoring}: come per il Processo di Gestione dei Cambiamenti
(Sezione~\ref{sec:cambiamenti}), lo standard ISO/IEC~12207:1995 non prevede un
processo di supporto autonomo per la gestione delle non conformità: la
individuazione, l'analisi e la risoluzione dei problemi (comprese le non
conformità) vi sono trattate nel Processo di Risoluzione dei Problemi
(\textit{Problem Resolution Process}, punto 6.8). Il gruppo ha scelto di
scorporarne la gestione in un processo a sé, estensione deliberata oltre lo
standard, per dare visibilità e tracciabilità autonome alle non conformità
rilevate durante verifica, review e retrospettiva, distinguendole dai
problemi tecnici generici.

\subsubsection{Descrizione}

Una non conformità è una deviazione rispetto alle presenti norme, al template, al
Piano di Qualifica o agli accordi formalizzati nei verbali.

Le non conformità possono emergere durante:

\begin{itemize}
    \item verifica tramite Pull Request;
    \item review di sprint;
    \item retrospettiva;
    \item revisione con committente;
    \item incontro con la proponente;
    \item controllo di un documento già rilasciato.
\end{itemize}

% ------------------------------------------------

\subsubsection{Classificazione}

Le non conformità vengono classificate in:

\begin{itemize}
    \item \textbf{Bloccante}: impedisce il merge, il rilascio o l'ingresso in baseline
    dell'artefatto;

    \item \textbf{Non Bloccante}: non impedisce il merge immediato, ma deve essere
    corretta tramite task tracciato nello sprint corrente o successivo.
\end{itemize}

% ------------------------------------------------

\subsubsection{Gestione}

Se la non conformità viene individuata prima del merge, il Verificatore la segnala
nella Pull Request tramite commento o richiesta di cambiamento.

Se la non conformità viene individuata dopo il rilascio, essa deve essere:

\begin{itemize}
    \item verbalizzata nella riunione successiva;
    \item trasformata in task;
    \item corretta tramite issue dedicata;
    \item verificata tramite Pull Request;
    \item registrata nel changelog del documento, se comporta modifica documentale.
\end{itemize}

Una non conformità bloccante può impedire il rilascio di una baseline fino alla
sua risoluzione.
